% journal:APEX
%DIF LATEXDIFF DIFFERENCE FILE
%DIF DEL 260222_MSEW_old.tex   Thu Mar 26 14:56:28 2026
%DIF ADD 260222_MSEW.tex       Thu Mar 26 12:22:12 2026
% author:Hiroki Yoshida
\documentclass[10pt]{iopart}
\usepackage{graphicx}
\usepackage{iopams}
\usepackage{bm}
\usepackage{hyperref}
%DIF PREAMBLE EXTENSION ADDED BY LATEXDIFF
%DIF UNDERLINE PREAMBLE %DIF PREAMBLE
\RequirePackage[normalem]{ulem} %DIF PREAMBLE
\RequirePackage{color}\definecolor{RED}{rgb}{1,0,0}\definecolor{BLUE}{rgb}{0,0,1} %DIF PREAMBLE
\providecommand{\DIFaddtex}[1]{{\protect\color{blue}\uwave{#1}}} %DIF PREAMBLE
\providecommand{\DIFdeltex}[1]{{\protect\color{red}\sout{#1}}}                      %DIF PREAMBLE
%DIF SAFE PREAMBLE %DIF PREAMBLE
\providecommand{\DIFaddbegin}{} %DIF PREAMBLE
\providecommand{\DIFaddend}{} %DIF PREAMBLE
\providecommand{\DIFdelbegin}{} %DIF PREAMBLE
\providecommand{\DIFdelend}{} %DIF PREAMBLE
\providecommand{\DIFmodbegin}{} %DIF PREAMBLE
\providecommand{\DIFmodend}{} %DIF PREAMBLE
%DIF FLOATSAFE PREAMBLE %DIF PREAMBLE
\providecommand{\DIFaddFL}[1]{\DIFadd{#1}} %DIF PREAMBLE
\providecommand{\DIFdelFL}[1]{\DIFdel{#1}} %DIF PREAMBLE
\providecommand{\DIFaddbeginFL}{} %DIF PREAMBLE
\providecommand{\DIFaddendFL}{} %DIF PREAMBLE
\providecommand{\DIFdelbeginFL}{} %DIF PREAMBLE
\providecommand{\DIFdelendFL}{} %DIF PREAMBLE
%DIF HYPERREF PREAMBLE %DIF PREAMBLE
\providecommand{\DIFadd}[1]{\texorpdfstring{\DIFaddtex{#1}}{#1}} %DIF PREAMBLE
\providecommand{\DIFdel}[1]{\texorpdfstring{\DIFdeltex{#1}}{}} %DIF PREAMBLE
\newcommand{\DIFscaledelfig}{0.5}
%DIF HIGHLIGHTGRAPHICS PREAMBLE %DIF PREAMBLE
\RequirePackage{settobox} %DIF PREAMBLE
\RequirePackage{letltxmacro} %DIF PREAMBLE
\newsavebox{\DIFdelgraphicsbox} %DIF PREAMBLE
\newlength{\DIFdelgraphicswidth} %DIF PREAMBLE
\newlength{\DIFdelgraphicsheight} %DIF PREAMBLE
% store original definition of \includegraphics %DIF PREAMBLE
\LetLtxMacro{\DIFOincludegraphics}{\includegraphics} %DIF PREAMBLE
\newcommand{\DIFaddincludegraphics}[2][]{{\color{blue}\fbox{\DIFOincludegraphics[#1]{#2}}}} %DIF PREAMBLE
\newcommand{\DIFdelincludegraphics}[2][]{% %DIF PREAMBLE
\sbox{\DIFdelgraphicsbox}{\DIFOincludegraphics[#1]{#2}}% %DIF PREAMBLE
\settoboxwidth{\DIFdelgraphicswidth}{\DIFdelgraphicsbox} %DIF PREAMBLE
\settoboxtotalheight{\DIFdelgraphicsheight}{\DIFdelgraphicsbox} %DIF PREAMBLE
\scalebox{\DIFscaledelfig}{% %DIF PREAMBLE
\parbox[b]{\DIFdelgraphicswidth}{\usebox{\DIFdelgraphicsbox}\\[-\baselineskip] \rule{\DIFdelgraphicswidth}{0em}}\llap{\resizebox{\DIFdelgraphicswidth}{\DIFdelgraphicsheight}{% %DIF PREAMBLE
\setlength{\unitlength}{\DIFdelgraphicswidth}% %DIF PREAMBLE
\begin{picture}(1,1)% %DIF PREAMBLE
\thicklines\linethickness{2pt} %DIF PREAMBLE
{\color[rgb]{1,0,0}\put(0,0){\framebox(1,1){}}}% %DIF PREAMBLE
{\color[rgb]{1,0,0}\put(0,0){\line( 1,1){1}}}% %DIF PREAMBLE
{\color[rgb]{1,0,0}\put(0,1){\line(1,-1){1}}}% %DIF PREAMBLE
\end{picture}% %DIF PREAMBLE
}\hspace*{3pt}}} %DIF PREAMBLE
} %DIF PREAMBLE
\LetLtxMacro{\DIFOaddbegin}{\DIFaddbegin} %DIF PREAMBLE
\LetLtxMacro{\DIFOaddend}{\DIFaddend} %DIF PREAMBLE
\LetLtxMacro{\DIFOdelbegin}{\DIFdelbegin} %DIF PREAMBLE
\LetLtxMacro{\DIFOdelend}{\DIFdelend} %DIF PREAMBLE
\DeclareRobustCommand{\DIFaddbegin}{\DIFOaddbegin \let\includegraphics\DIFaddincludegraphics} %DIF PREAMBLE
\DeclareRobustCommand{\DIFaddend}{\DIFOaddend \let\includegraphics\DIFOincludegraphics} %DIF PREAMBLE
\DeclareRobustCommand{\DIFdelbegin}{\DIFOdelbegin \let\includegraphics\DIFdelincludegraphics} %DIF PREAMBLE
\DeclareRobustCommand{\DIFdelend}{\DIFOaddend \let\includegraphics\DIFOincludegraphics} %DIF PREAMBLE
\LetLtxMacro{\DIFOaddbeginFL}{\DIFaddbeginFL} %DIF PREAMBLE
\LetLtxMacro{\DIFOaddendFL}{\DIFaddendFL} %DIF PREAMBLE
\LetLtxMacro{\DIFOdelbeginFL}{\DIFdelbeginFL} %DIF PREAMBLE
\LetLtxMacro{\DIFOdelendFL}{\DIFdelendFL} %DIF PREAMBLE
\DeclareRobustCommand{\DIFaddbeginFL}{\DIFOaddbeginFL \let\includegraphics\DIFaddincludegraphics} %DIF PREAMBLE
\DeclareRobustCommand{\DIFaddendFL}{\DIFOaddendFL \let\includegraphics\DIFOincludegraphics} %DIF PREAMBLE
\DeclareRobustCommand{\DIFdelbeginFL}{\DIFOdelbeginFL \let\includegraphics\DIFdelincludegraphics} %DIF PREAMBLE
\DeclareRobustCommand{\DIFdelendFL}{\DIFOaddendFL \let\includegraphics\DIFOincludegraphics} %DIF PREAMBLE
%DIF COLORLISTINGS PREAMBLE %DIF PREAMBLE
\RequirePackage{listings} %DIF PREAMBLE
\RequirePackage{color} %DIF PREAMBLE
\lstdefinelanguage{DIFcode}{ %DIF PREAMBLE
%DIF DIFCODE_UNDERLINE %DIF PREAMBLE
  moredelim=[il][\color{red}\sout]{\%DIF\ <\ }, %DIF PREAMBLE
  moredelim=[il][\color{blue}\uwave]{\%DIF\ >\ } %DIF PREAMBLE
} %DIF PREAMBLE
\lstdefinestyle{DIFverbatimstyle}{ %DIF PREAMBLE
	language=DIFcode, %DIF PREAMBLE
	basicstyle=\ttfamily, %DIF PREAMBLE
	columns=fullflexible, %DIF PREAMBLE
	keepspaces=true %DIF PREAMBLE
} %DIF PREAMBLE
\lstnewenvironment{DIFverbatim}{\lstset{style=DIFverbatimstyle}}{} %DIF PREAMBLE
\lstnewenvironment{DIFverbatim*}{\lstset{style=DIFverbatimstyle,showspaces=true}}{} %DIF PREAMBLE
%DIF END PREAMBLE EXTENSION ADDED BY LATEXDIFF

\begin{document}
\title{Magnetoelastic wave in Ferromagnetic Thin-film Mediated by Dipolar Interaction}
% \article{\APEX}{Magnetoelastic wave in Ferromagnetic Thin-film Mediated by Dipolar Interaction}

\author{Hiroki Yoshida$^1$, Ryohei Kono$^1$, Manato Fujimoto$^2$, Motoki Asano$^3$, Daiki Hatanaka$^3$, Kei Yamamoto$^{4,5}$, Shuichi Murakami$^{2,5,6}$}
\address{$^1$ Department of Physics, Institute of Science Tokyo, 2-12-1 Ookayama, Meguro-ku, Tokyo 152-8551, Japan}
\address{$^2$ Department of Applied Physics, The University of Tokyo, 7-3-1 Hongo, Bunkyo-ku, Tokyo 113-8656, Japan}
\address{$^3$ Basic Research Laboratories, NTT, Inc., 3-1 Morinosato-Wakamiya, Atsugi-shi, Kanagawa 243-0198, Japan}
\address{$^4$ Advanced Science Research Center, Japan Atomic Energy Agency, Tokai, Ibaraki 319-1195, Japan}
\address{$^5$ Center for Emergent Matter Science, RIKEN, 2-1 Hirosawa, Wako, Saitama 351-0198, Japan}
\address{$^6$ International Institute for Sustainability with Knotted Chiral Meta Matter (WPI-SKCM$^2$), Hiroshima University, Higashi-hiroshima, Hiroshima 739-0046, Japan}
\ead{yoshida.h.9d8d@m.isct.ac.jp}


\begin{abstract}
    Magnetoelastic coupling mediated by magnetic dipolar \DIFdelbegin \DIFdel{interactions }\DIFdelend \DIFaddbegin \DIFadd{interaction }\DIFaddend is theoretically investigated in \DIFdelbegin \DIFdel{ferromagnetic thin films }\DIFdelend \DIFaddbegin \DIFadd{a ferromagnetic thin film }\DIFaddend under an in-plane magnetic field. We \DIFdelbegin \DIFdel{construct a microscopic framework }\DIFdelend \DIFaddbegin \DIFadd{develop a theoretical description }\DIFaddend that incorporates dipolar fields derived from Maxwell's equations in the presence of elastic deformation. The resulting coupled equations of motion predict hybridization between magnetostatic and Lamb waves. Numerical calculations for a yttrium iron garnet (YIG) film reveal \DIFdelbegin \DIFdel{hybridization of modes and }\DIFdelend anti-crossings \DIFdelbegin \DIFdel{of $0.1$ to several MHz }\DIFdelend in the dispersion relations\DIFaddbegin \DIFadd{, with hybridization gaps ranging from $0.1$ to several MHz}\DIFaddend .
\end{abstract}

% \keywords{Universe}
\submitto{\APEX}
\maketitle
\ioptwocol

Magnetoelastic coupling has attracted considerable attention because of its potential for manipulating magnetic properties via mechanical degrees of freedom~\cite{Kittel_PR,LL_EDinCM}. At the microscopic level, magnons and phonons are coupled \DIFdelbegin \DIFdel{through }\DIFdelend \DIFaddbegin \DIFadd{via }\DIFaddend spin-orbit interaction\DIFdelbegin \DIFdel{, which gives rise to magnon-phonon coupling}\DIFdelend ~\cite{Kittel_PR,Kittel_Abraham_RMP,Flebus_magnonPolaron}. In contrast, at longer wavelengths\DIFdelbegin \DIFdel{the coupling is mainly mediated by the magnetic dipolar interaction. In this regime, magnetostatic waves interact with elastic waves. Experimental studies of such dipolar-interaction-mediated magnetoelastic coupling have recently progressed owing to the development of }\DIFdelend \DIFaddbegin \DIFadd{, magnetic dipolar interactions also play an important role for the coupling. Recent experimental advances, particularly those based on }\DIFaddend interdigital transducers (IDTs)\DIFdelbegin \DIFdel{~\mbox{%DIFAUXCMD
\cite{IDT_review,Kunstle2025}}\hskip0pt%DIFAUXCMD
, which enable precise measurements of acoustic excitations }\DIFdelend \DIFaddbegin \DIFadd{, have enabled coherent excitation and precise characterization of acoustic waves }\DIFaddend in magnetic thin films\DIFdelbegin \DIFdel{. However, a }\DIFdelend \DIFaddbegin \DIFadd{~\mbox{%DIFAUXCMD
\cite{IDT_review,Kunstle2025}}\hskip0pt%DIFAUXCMD
, allowing direct access to magnetoelastic coupling in this regime. While continuum theories describing magnetostatic~\mbox{%DIFAUXCMD
\cite{Walker1957,damoneshbachMagnetostaticModesFerromagnet1961,damonPropagationMagnetostaticSpin} }\hskip0pt%DIFAUXCMD
and elastic waves~\mbox{%DIFAUXCMD
\cite{Rayleigh1885,Rayleigh1889,Lamb1889,lambWavesElasticPlate1917} }\hskip0pt%DIFAUXCMD
in this regime have been developed separately, a unified }\DIFaddend theoretical framework that \DIFdelbegin \DIFdel{consistently incorporates the dipolar interaction in magnetoelastic dynamics has not yet been established, because of its long-ranged feature and its singularity when two dipoles comes close}\DIFdelend \DIFaddbegin \DIFadd{captures their coupling mediated by dipolar interactions remains lacking, primarily due to the long-range nature of the dipolar interaction and its singular behavior at short distances}\DIFaddend .

In this Letter, \DIFdelbegin \DIFdel{by circumventing the above difficulties using the Green function method, we construct a general formalism that incorporates the dipolar interaction and describe the resulting coupling between magnetostatic and elastic waves in ferromagnetic thin films. }\DIFdelend \DIFaddbegin \DIFadd{we extend the Green-function method developed by Kalinikos and Slavin~\mbox{%DIFAUXCMD
\cite{Kalinikos_Slavin_1986} }\hskip0pt%DIFAUXCMD
to incorporate dipolar interactions in the presence of elastic deformations. This approach allows us to quantitatively evaluate the magnetoelastic coupling mediated by dipolar interactions.
}\DIFaddend 

Our system is schematically shown in Fig.~\ref{fig:Schematic}. We consider a ferromagnetic thin film that is infinitely extended in the $xy$ plane and has a finite thickness $L$ in the $z$ direction, with the origin taken at the middle of the film. \DIFdelbegin \DIFdel{The film }\DIFdelend \DIFaddbegin \DIFadd{It }\DIFaddend is assumed to be isotropic both magnetically and elastically. An in-plane external magnetic field $\bm{H}_0$ is applied, resulting in an equilibrium saturation magnetization $\bm{M}_0$ parallel to $\bm{H}_0$.

\begin{figure}
    \begin{center}
        \includegraphics[width=\columnwidth]{fig/Schematic_MSEW.pdf}
        \caption{Schematic illustration of the system under consideration. A ferromagnetic thin film with thickness $L$ in the $z$ direction and infinite extent in the $x$ and $y$ directions is considered. Magnetic moments precess around an axis parallel to the in-plane external magnetic field $\bm{H}_0$. The elastic wave changes the positions of the magnetic moments and their separations, resulting in a modification of the dipolar interaction among them.}
        \label{fig:Schematic}
    \end{center}
\end{figure}

We decompose the magnetic field $\bm{H}$ and the magnetization $\bm{M}$ as $\bm{H}=\bm{H}_0+\bm{h}(\bm{r},t)$ and $\bm{M}=\bm{M}_0+\bm{m}(\bm{r},t)$, respectively\DIFaddbegin \DIFadd{, where $\bm{r}=(x,y,z)$ is the Euler coordinate, }\DIFaddend and study propagating magnetostatic waves of small deviations from the equilibrium magnetization $\bm{m}(\bm{r},t)$. The system also supports elastic waves described by the displacement $\bm{u}(\bm{r},t)$. The deformation modulates the distance between magnetic dipoles, so that the \DIFdelbegin \DIFdel{dipole }\DIFdelend \DIFaddbegin \DIFadd{dipolar }\DIFaddend field $\bm{h}(\bm{r},t)$ generated by $\bm{m}$ is influenced by $\bm{u}$, leading to a magnetoelastic wave. We take the $x$ axis parallel to the propagation direction, and assume that waves are plane waves and the system is uniform in the $y$ direction\DIFdelbegin \DIFdel{, where we write $\bm{r}=(x,y,z)$}\DIFdelend . Then, $\bm{m}$ and $\bm{u}$ can be written as
\begin{equation}
    \bm{m}(\bm{r},t)=\bm{m}(z)e^{i(kx-\omega t)},\\
    \bm{u}(\bm{r},t)=\bm{u}(z)e^{i(kx-\omega t)}\DIFdelbegin \DIFdel{.}\DIFdelend \DIFaddbegin \DIFadd{,}\DIFaddend \label{eq:Assumption}
\end{equation}
\DIFaddbegin \DIFadd{where $\bm{m}(z),\bm{u}(z)\in\mathbb{C}^3$. Hereafter, all fields are treated as complex quantities, and physical observables are given by their real parts. }\DIFaddend We study the coupled dynamics of these waves up to first order in these quantities. Here, we focus exclusively on the case of an in-plane magnetic field. The case of an arbitrary external magnetic field will be presented elsewhere.

We start from the basic equations describing the magnetic dipolar interaction, namely Maxwell's equations. When elastic deformation is present, the magnetic dipoles deviate from their equilibrium positions $\bm{r}$ to $\bm{R}:=\bm{r}+\bm{u}(\bm{r},t)$. The Maxwell's equations are then rewritten in terms of this \DIFdelbegin \DIFdel{new }\DIFdelend \DIFaddbegin \DIFadd{Lagrange }\DIFaddend coordinate as
\begin{equation}
    \nabla_{\bm{R}}\cdot \left(\bm{H}+\bm{M}\right)=0,\qquad
    \nabla_{\bm{R}}\times \bm{H}\approx 0,\label{eq:Maxwell_R}
\end{equation}
where in the second equation the magnetostatic approximation is applied and the time derivative of the electric field is discarded. This approximation is justified when the \DIFdelbegin \DIFdel{length scale is sufficiently long so that the dipolar interaction is dominant}\DIFdelend \DIFaddbegin \DIFadd{steady current is absent and the characteristic frequency $\omega$ and length scale $L$ of the magnetization dynamics satisfy the quasi-static condition $\omega L/c \ll 1$, so that retardation effects and electromagnetic radiation can be ignored~\mbox{%DIFAUXCMD
\cite{Stancil_TheoryofMagnetostaticWaves}}\hskip0pt%DIFAUXCMD
}\DIFaddend . By the plane-wave assumption~(\ref{eq:Assumption}), the dipolar field can also be written as $\bm{h}=\bm{h}(z)e^{i(kx-\omega t)}$. We rewrite Eq.~(\ref{eq:Maxwell_R}) in terms of these quantities in the \DIFdelbegin \DIFdel{original }\DIFdelend \DIFaddbegin \DIFadd{Euler }\DIFaddend coordinate $\bm{r}$. We note that we need to pay attention to the conservation of the magnetization in a unit volume~\cite{Yamamoto_Maekawa_AnnPhys} yielding $\bm{M}(\bm{R})=\bm{M}(\bm{r})/\det(J(\bm{r}))$, where $J(\bm{r})$ is a Jacobian matrix (see Supplementary data~\cite{SM} for details). We get
\begin{equation}
    \nabla_{\bm{r}}\cdot\left(\bm{h}(\bm{r})+\bm{m}(\bm{r})\right)=\bm{M}_0\cdot\nabla_{\bm{r}}\left(\nabla_{\bm{r}}\cdot\bm{u}(\bm{r})\right),\label{eq:Maxwell_couple_div}
\end{equation}
\begin{equation}
    \nabla_{\bm{r}}\times\bm{h}(\bm{r})=0,\label{eq:Maxwell_couple_rot}
\end{equation}
where $\nabla_{\bm{r}}$ stands for derivatives with respect to the $\bm{r}$ coordinate. Because the surface of the film is also distorted as in Fig.~\ref{fig:Schematic}, the boundary conditions used to solve these equations are also modified by the presence of $\bm{u}$ as
\begin{equation}
    h^{\mathrm{in}}_x-h_x^{\mathrm{out}}=\DIFdelbegin \DIFdel{M_{0,z}\frac{\partial u_z}{\partial x}}\DIFdelend \DIFaddbegin \DIFadd{0}\DIFaddend ,\label{eq:Maxwell_bc_x}
\end{equation}
\begin{equation}
    h^{\mathrm{in}}_z+m_z-h^{\mathrm{out}}_z=M_{0,x}\frac{\partial u_z}{\partial x},\label{eq:Maxwell_bc_z}
\end{equation}
at $z=\pm L/2$, where $h^{\mathrm{in}/\mathrm{out}}$ stands for the dipolar fields inside and outside the film, respectively. Hereafter, lower indices indicate \DIFdelbegin \DIFdel{component }\DIFdelend \DIFaddbegin \DIFadd{components }\DIFaddend in the $\bm{r}$ coordinate. Following the Green tensor method as in Ref.~\cite{Kalinikos_Slavin_1986}, the dipolar field satisfying Eqs.~(\ref{eq:Maxwell_couple_div}) and (\ref{eq:Maxwell_couple_rot}) under the boundary conditions (\ref{eq:Maxwell_bc_x}) and (\ref{eq:Maxwell_bc_z}) is given as (see Supplementary data~\cite{SM} for details)
\begin{equation}
    \bm{h}(z)=\int_{-\frac{L}{2}}^{\frac{L}{2}}\rmd z'\left(G^m(z,z')\bm{m}(z')+G^u(z,z')\bm{u}(z')\right),\label{eq:dipolar_field}
\end{equation}
where
\begin{equation}
    G^m(z,z')=\left(\begin{array}{ccc}
        -G_P(z,z')&0&-iG_Q(z,z')\\
        0&0&0\\
        -iG_Q(z,z')&0&G_{P'}(z,z')
    \end{array}\right),
\end{equation}
\begin{equation*}
    G^u(z,z')=-ikM_{0,x}G^m(z,z'),
\end{equation*}
\begin{equation}
    G_{P}(z,z')=\frac{k}{2}e^{-k|z-z'|},
\end{equation}
\begin{equation}
    G_{Q}(z,z')=\mathrm{sgn}(z-z')G_{P}(z,z'),
\end{equation}
\begin{equation}
    G_{P'}(z,z')=G_{P}(z,z')-\delta(z-z').
\end{equation}
This dipolar field couples the magnetostatic wave and the elastic wave in the film. We note that in the present setup there is no coupling when the equilibrium magnetization is parallel to the $y$ axis\DIFdelbegin \DIFdel{, i.e., when the external magnetic field is applied in that direction.
Otherwise, the two waves are coupled via $\bm{h}$}\DIFdelend .

Next, we examine the effect of magnetoelastic coupling on the dynamics of the magnetostatic and elastic waves. In the absence of this coupling, the dynamics of the magnetization $\bm{m}$ up to first order is described by the linearized Landau--Lifshitz equation
\begin{equation}
    \dot{\bm{m}}=-\gamma\mu_0\left(\bm{m}\times\bm{H}_0+\bm{M}_0\times\bm{h}\right),\label{eq:LandauLifshitz}
\end{equation}
where $\gamma$ is the gyromagnetic ratio and $\mu_0$ is the vacuum magnetic permeability. This equation can be derived from a Lagrangian
\begin{equation*}
    L^m[\bm{m},\dot{\bm{m}},\bm{m}^*,\dot{\bm{m}}^*]=\int\rmd z\Re\left[\frac{1}{4M_0^2\gamma}\bm{M}_0\cdot(\bm{m}\times\dot{\bm{m}}^*)\right.
\end{equation*}
\begin{equation}
    \hspace{2.9cm}\left.-\frac{\mu_0}{4}\left(\bm{h}^*\cdot\bm{h}+\frac{H_0}{M_0}\bm{m}^*\cdot\bm{m}\right)\right]\DIFdelbegin \DIFdel{.}\DIFdelend \DIFaddbegin \DIFadd{,}\DIFaddend \label{eq:Lagrangian_m}
\end{equation}
\DIFaddbegin \DIFadd{where the asterisk shows complex conjugate. }\DIFaddend We can verify this by constructing the Euler--Lagrange equation with respect to $\bm{m}^*$ from this Lagrangian. On the other hand, the dynamics of the displacement field $\bm{u}$ is described by the Navier--Cauchy equation
\begin{equation}
    \rho\ddot{\bm{u}}=\mu\nabla^2\bm{u}+\left(\mu+\lambda\right)\nabla\left(\nabla\cdot\bm{u}\right),\label{eq:NavierCauchy}
\end{equation}
where $\rho$ is the mass density, $\mu$ and $\lambda$ are Lam\'e constants. The Lagrangian for this equation is
\begin{equation}
    L^u[\bm{u},\dot{\bm{u}},\bm{u}^*,\dot{\bm{u}}^*]=\int\rmd z\Re\left[\frac{\rho}{4}\dot{\bm{u}}^*\cdot\dot{\bm{u}}+\frac{1}{4}\partial_i\sigma_{ij}u_j^*\right],\label{eq:Lagrangian_u}
\end{equation}
where $\sigma_{ij}=\mu(\partial_i u_j+\partial_j u_i)+\lambda\delta_{ij}\nabla\cdot\bm{u}$ is the stress tensor. The total Lagrangian of the coupled system is given by the sum of the two Lagrangians, Eqs.~(\ref{eq:Lagrangian_m}) and (\ref{eq:Lagrangian_u}). In this case, $\bm{h}$ is no longer only a functional of the magnetization, but also a functional of the displacement field $\bm{u}$. As a result, a new volume force term \DIFdelbegin \DIFdel{$-\mu_0\bm{f}$ }\DIFdelend \DIFaddbegin \DIFadd{$-\bm{f}$ }\DIFaddend appears in the Navier--Cauchy equation~(\ref{eq:NavierCauchy}), where the force \DIFaddbegin \DIFadd{density }\DIFaddend $\bm{f}$ is given by (see Supplementary data~\cite{SM} for details of the derivation)
\begin{equation}
    f_i(z):=\DIFaddbegin \DIFadd{\mu_0}\DIFaddend \frac{\delta h^*_j}{\delta u^*_i}h_j=\DIFdelbegin \DIFdel{-ikM}\DIFdelend \DIFaddbegin \DIFadd{-ik\mu_0 M}\DIFaddend _{0,x}h_i(z).
\end{equation}
Physically, this force can be understood as a Maxwell stress acting on the magnetic dipoles, caused by the change in the dipolar field induced by the displacement field. The coupled dynamics of the magnetoelastic wave is then given by
\begin{equation}
    \left\{
    \DIFdelbegin %DIFDELCMD < \begin{array}{l}
%DIFDELCMD <         \dot{\bm{m}}=-\gamma\mu_0\left(\bm{m}\times\bm{H}_0+\bm{M}_0\times\bm{h}\right)\\
%DIFDELCMD <         \rho\ddot{\bm{u}}=\mu\nabla^2\bm{u}+\left(\mu+\lambda\right)\nabla\left(\nabla\cdot\bm{u}\right)-\mu_0\bm{f}.
%DIFDELCMD <     \end{array}%%%
\DIFdelend \DIFaddbegin \begin{array}{l}
        \dot{\bm{m}}=-\gamma\mu_0\left(\bm{m}\times\bm{H}_0+\bm{M}_0\times\bm{h}\right)\\
        \rho\ddot{\bm{u}}=\mu\nabla^2\bm{u}+\left(\mu+\lambda\right)\nabla\left(\nabla\cdot\bm{u}\right)-\bm{f}.
    \end{array}\DIFaddend \label{eq:CoupledEOM}
    \right.
\end{equation}


\begin{figure*}[t]
    \begin{center}
        \includegraphics[width=2 \columnwidth]{fig/Bands.pdf}
        \caption{Dispersion curves of the coupled waves. Red and blue curves represent symmetric and antisymmetric modes of the uncoupled waves, including magnetostatic backward volume waves (BVWs) and Lamb waves, while black curves represent magnetostatic modes that are neither symmetric nor antisymmetric. \DIFdelbeginFL \DIFdelFL{Black dots }\DIFdelendFL \DIFaddbeginFL \DIFaddFL{Markers }\DIFaddendFL show the dispersion of the coupled waves obtained from numerical calculations. The external magnetic field is applied (a) parallel to the propagation direction, $\bm{H}_0 \parallel \bm{k}$ ($\varphi=0$), (b) at $\varphi=\pi/4$, and (c) perpendicular to the propagation direction, $\bm{H}_0 \perp \bm{k}$ ($\varphi=\pi/2$). Insets show magnified views of representative band crossings. In (a) $\varphi=0$, anticrossings appear between BVW and Lamb modes only when they have the same $C_{2x}$ eigenvalues (same color). When the $C_{2x}$ symmetry of the system is broken (b), both the magnetostatic surface mode and the BVW couple to Lamb modes. \DIFaddbeginFL \DIFaddFL{The bottom-center inset of panel (b) shows the hybridization between the lowest-frequency backward volume wave mode and the first Lamb wave mode.}\DIFaddendFL }
        \label{fig:Bands}
    \end{center}
\end{figure*}

In order to examine the dynamical properties of the coupled waves described by the above equations of motion, we numerically calculate the dispersion relation in the following. Under the assumption~(\ref{eq:Assumption}), the time \DIFdelbegin \DIFdel{derivative }\DIFdelend \DIFaddbegin \DIFadd{derivatives }\DIFaddend on the left-hand side of the equations \DIFdelbegin \DIFdel{becomes }\DIFdelend \DIFaddbegin \DIFadd{become }\DIFaddend $-i\omega$. The second \DIFdelbegin \DIFdel{equation }\DIFdelend of Eqs.~(\ref{eq:CoupledEOM}) contains spatial derivatives of $\bm{u}$ and requires boundary conditions to \DIFdelbegin \DIFdel{solve it}\DIFdelend \DIFaddbegin \DIFadd{be solved}\DIFaddend . Since we consider a film in vacuum, we impose the \DIFdelbegin \DIFdel{tension-free }\DIFdelend \DIFaddbegin \DIFadd{stress-free }\DIFaddend condition at the surfaces $z=\pm L/2$ as
\begin{equation}
    \sigma_{xz}|_{z=\pm\frac{L}{2}}=\sigma_{yz}|_{z=\pm\frac{L}{2}}=\sigma_{zz}|_{z=\pm\frac{L}{2}}=0.\label{eq:BC_EW}
\end{equation}
Without the coupling, it is known that a type of elastic wave called a Lamb wave appears under these conditions. Since the system in equilibrium is symmetric with respect to the $z=0$ plane, the Lamb wave modes are either \DIFdelbegin \DIFdel{symmetric or antisymmetric with respect to the sign change of the variable }\DIFdelend \DIFaddbegin \DIFadd{even or odd in }\DIFaddend $z$. Numerically, we need to solve the equations under the boundary conditions~(\ref{eq:BC_EW}). For example, the finite element method naturally incorporates these conditions.

On the other hand, the magnetostatic waves without the coupling to elastic waves in this geometry were studied in detail by Damon and Eshbach~\DIFdelbegin \DIFdel{\mbox{%DIFAUXCMD
\cite{damonMagnetostaticModesFerromagnet1961}}\hskip0pt%DIFAUXCMD
}\DIFdelend \DIFaddbegin \DIFadd{\mbox{%DIFAUXCMD
\cite{damoneshbachMagnetostaticModesFerromagnet1961}}\hskip0pt%DIFAUXCMD
}\DIFaddend . They analytically showed that when the external magnetic field is parallel to the propagation direction of the wave ($\bm{H}_0\parallel\bm{k}$), the system has infinitely many bulk modes called backward volume waves (BVWs), which are either symmetric or antisymmetric. When the magnetic field is in-plane but perpendicular to the propagation direction ($\bm{H}_0\perp\bm{k}$), a surface mode appears that is localized at one of the surfaces.

For the calculation of the dispersion, it is convenient to introduce another coordinate \DIFdelbegin \DIFdel{$(r_1,r_2,r_3)$}\DIFdelend \DIFaddbegin \DIFadd{system}\DIFaddend , where the original \DIFdelbegin \DIFdel{coordinate }\DIFdelend \DIFaddbegin \DIFadd{one }\DIFaddend $(x,y,z)$ is rotated around the $z$ axis so that \DIFdelbegin \DIFdel{the $r_1$ }\DIFdelend \DIFaddbegin \DIFadd{one of the }\DIFaddend axis is parallel to the external magnetic field. Then, the \DIFdelbegin \DIFdel{$r_1$ }\DIFdelend component of $\bm{m}$ is zero \DIFaddbegin \DIFadd{along this direction }\DIFaddend since $\bm{M}_0\cdot\bm{m}=0$ from the constraint $|\bm{M}|=const.$
\DIFdelbegin \DIFdel{.
}\DIFdelend 

Now we consider the coupled equations~(\ref{eq:CoupledEOM}). Using the finite element method, we discretize the equations in real space while satisfying the boundary conditions, which leads to a generalized eigenvalue problem for the eigenvalue $\omega$. By solving this problem, we obtain the dispersion curves shown in Fig.~\ref{fig:Bands}. The parameters are taken to be those of Yttrium-Iron-Garnet (YIG): $\rho=5.17\ \mathrm{g}/\mathrm{cm}^3$, $\mu=7.81\times10^{-8}\ \mathrm{kg}\cdot\mathrm{GHz}^2/\mathrm{m}$, $\lambda=10.94\times10^{-8}\ \mathrm{kg}\cdot\mathrm{GHz}^2/\mathrm{m}$~\cite{clarkElasticConstantsSingleCrystal1961}, $\gamma=29.8\ \mathrm{GHz}/\mathrm{T}$, and $\mu_0M_0=0.178$ T~\cite{AulockHandbookofferritematerials}. The vacuum permeability is $\mu_0=4\pi\times10^{11}\ \mathrm{T}^2\mathrm{m}/(\mathrm{kg}\cdot\mathrm{GHz}^2)$. We set the film thickness to \DIFdelbegin \DIFdel{$L=500\ \mu$m }\DIFdelend \DIFaddbegin \DIFadd{$L=500$ nm }\DIFaddend and the external magnetic field to $\mu_0H_0=0.1$ T.

We define the angle between the $x$ axis and the external magnetic field as $\varphi$ and write $\bm{H}_0=H_0(\cos\varphi,\sin\varphi,0)$. In Fig.~\ref{fig:Bands}(a), we plot the case $\varphi=0$. In this case, \DIFdelbegin \DIFdel{magnetostatic BVWs appear, plotted as blue and red curves representing the antisymmetric and symmetric modes with respect to the twofold rotation about the $x$-axis (}\DIFdelend \DIFaddbegin \DIFadd{the system has a }\DIFaddend $C_{2x}$ \DIFdelbegin \DIFdel{), with eigenvalues $-1$ }\DIFdelend \DIFaddbegin \DIFadd{rotation symmetry and thus, the modes of magnetostatic BVWs }\DIFaddend and \DIFaddbegin \DIFadd{Lamb waves are eigenstates of the $C_{2x}$ operation. States with $C_{2x}$ eigenvalue $-1$ ($C_{2x}\bm{v}(z)=-\bm{v}(z),\,(\bm{v}=\bm{m},\bm{u})$) are plotted as blue curves, while those with }\DIFaddend $+1$ \DIFdelbegin \DIFdel{, respectively. Similarly, antisymmetric and symmetric Lamb waves with positive group velocities (positive slopes ) are shown by blue and red curves}\DIFdelend \DIFaddbegin \DIFadd{as red curves. Curves with negative slopes are magnetostatic BVWs and positive ones are Lamb waves}\DIFaddend , respectively. The \DIFdelbegin \DIFdel{dots }\DIFdelend \DIFaddbegin \DIFadd{markers }\DIFaddend represent the dispersion of the coupled waves. The insets show magnified views of the band crossings between \DIFdelbegin \DIFdel{BVW }\DIFdelend \DIFaddbegin \DIFadd{BVWs }\DIFaddend and Lamb waves. We observe that the magnetostatic and elastic waves \DIFdelbegin \DIFdel{couple and produce anti-crossings }\DIFdelend \DIFaddbegin \DIFadd{cause anti-crossing and produce hybridization gaps }\DIFaddend of several MHz only when the two modes have the same $C_{2x}$ eigenvalues. \DIFdelbegin \DIFdel{This is because when $\bm{H}_0\parallel\bm{k}$, the system possesses the two-fold rotation symmetry with respect to the $x$ axis $C_{2x}$ and coupled modes needs to be either symmetric or anti-symmetric. }\DIFdelend Figure~\ref{fig:Bands}(b) shows the case $\varphi=\pi/4$. In this case, two types of magnetostatic modes coexist: BVWs and surface \DIFdelbegin \DIFdel{modes and the $C_{2x}$ symmetry is broken}\DIFdelend \DIFaddbegin \DIFadd{waves}\DIFaddend . The insets show magnified views of the band crossings between the magnetostatic and Lamb \DIFdelbegin \DIFdel{modes. Small }\DIFdelend \DIFaddbegin \DIFadd{waves. Hybridization gaps at }\DIFaddend anti-crossings appear between all the modes as the \DIFdelbegin \DIFdel{two-fold rotation symmetry is broken }\DIFdelend \DIFaddbegin \DIFadd{system does not have any symmetries }\DIFaddend in this case. Finally, in Fig.~\ref{fig:Bands}(c), we plot the case $\varphi=\pi/2$. \DIFdelbegin \DIFdel{In this case}\DIFdelend \DIFaddbegin \DIFadd{Here}\DIFaddend , a magnetostatic surface mode appears but does not couple to the Lamb \DIFdelbegin \DIFdel{modes}\DIFdelend \DIFaddbegin \DIFadd{waves}\DIFaddend . This is because $G^u=0,\bm{f}=\bm{0}$ and Eq.~(\ref{eq:CoupledEOM}) reduces to two independent equations. In other words, the magnetostatic modes are odd under the mirror reflection with respect to the $xz$ plane, while the elastic modes are even, and therefore they do not couple.

\DIFaddbegin \DIFadd{From the magnitude of the hybridization gaps, the effective coupling strength $g_{\mathrm{eff}}$ is estimated to be on the order of a few MHz. On the other hand, assuming a typical Gilbert damping parameter $\alpha \approx 10^{-4}$~\mbox{%DIFAUXCMD
\cite{YIGalpha_PRAppl}}\hskip0pt%DIFAUXCMD
, a quality factor $Q \approx 10^3$, and operating frequencies $\omega$ in the GHz range, the corresponding magnon and phonon linewidths are estimated as $\gamma_m \sim \alpha \omega$ and $\kappa \sim \omega/Q$, respectively. Under these conditions, $g_{\mathrm{eff}}^2/(\gamma_m\kappa)\sim1$, suggesting that the system is close to the boundary of the strong-coupling regime~\mbox{%DIFAUXCMD
\cite{rodriguezClassicalQuantumDistinctions2016}}\hskip0pt%DIFAUXCMD
. Therefore, the hybridization gaps are expected to be experimentally resolvable.
}

\DIFaddend In conclusion, we have \DIFdelbegin \DIFdel{constructed a theoretical framework to describe }\DIFdelend \DIFaddbegin \DIFadd{developed a theoretical description of }\DIFaddend magnetoelastic coupling mediated by magnetic dipolar interactions in ferromagnetic thin films. The dipolar field was derived from Maxwell's equations in the presence of elastic deformations, and the coupled equations of motion were obtained within the Lagrangian formalism. Numerical calculations were performed for a YIG thin film to obtain the dispersion relations of the coupled modes. We found \DIFdelbegin \DIFdel{mode hybridizations and resulting }\DIFdelend anti-crossings \DIFaddbegin \DIFadd{due to mode hybridization, with gaps }\DIFaddend of several MHz in general. When the external in-plane magnetic field is parallel to the \DIFdelbegin \DIFdel{propagation of the waves, }\DIFdelend \DIFaddbegin \DIFadd{wave propagation, the system has $C_2$ symmetry, so }\DIFaddend only magnetostatic and elastic waves with the same \DIFdelbegin \DIFdel{parity }\DIFdelend \DIFaddbegin \DIFadd{$C_2$ symmetry eigenvalues }\DIFaddend can hybridize. When they are perpendicular, couplings \DIFdelbegin \DIFdel{does }\DIFdelend \DIFaddbegin \DIFadd{do }\DIFaddend not occur. The formulation developed here provides a starting point for theoretical studies of dipolar-interaction-mediated magnetoelastic coupling in ferromagnetic thin films.

\ack{H. Y. acknowledges support from Japan Society for the Promotion of Science (JSPS) KAKENHI Grant No. JP24KJ1109 and by MEXT Initiative to Establish Next-generation Novel Integrated Circuits Centers (X-NICS) Grant No. JPJ011438. S. M. acknowledges support by Japan Society for the Promotion of Science (JSPS) KAKENHI Grant No. JP22K18687, No. JP22H00108, and No. JP24H02231.}

\bibliographystyle{iopart-num.bst}
\bibliography{Ref_MSEW.bib}

\end{document}